% Independent mathematical notes for integration into the main report.
\section{An independent lower bound from a helix}
Let
\[
 r=\frac1{\pi\sqrt6},\qquad h=\frac1{\sqrt3},\qquad
 \gamma(t)=\left(r\cos t,r\sin t,\frac{ht}{2\pi}\right),
 \quad 0\le t\le2\pi.
\]
This is a smooth, simple, one-turn circular helix. Its length is
\[
 2\pi\sqrt{r^2+\left(\frac h{2\pi}\right)^2}
 =\sqrt{4\pi^2r^2+h^2}=1.
\]
We prove directly that any convex body containing this curve has volume at
least $\pi r^2h/3$, without relying on a formula for the entire convex hull.

Write $A=\gamma(0)$ and $B=\gamma(2\pi)$, and let $p=(r,0)$ be their common
orthogonal projection onto the $xy$-plane. For $N\ge3$, set
$q_i=\gamma(2\pi i/N)$, $0\le i\le N$. Every tetrahedron
\[
 T_i=\operatorname{conv}\{A,B,q_i,q_{i+1}\},\qquad 1\le i\le N-2,
\]
lies in the convex container. Its orthogonal projection is the triangle
with vertices $p$, $\pi_{xy}(q_i)$, $\pi_{xy}(q_{i+1})$. These triangles form
the usual fan triangulation of the regular inscribed $N$-gon. In particular,
their interiors are pairwise disjoint; the tetrahedra therefore also have
pairwise disjoint interiors. Since $B-A=(0,0,h)$, the determinant formula for
tetrahedron volume gives
\[
 |T_i|=\frac h3\operatorname{area}
 \operatorname{conv}\{p,\pi_{xy}(q_i),\pi_{xy}(q_{i+1})\}.
\]
Adding yields
\[
 |K|\ge\frac h3\frac{Nr^2}{2}\sin\frac{2\pi}{N}.
\]
Letting $N\to\infty$ proves
\[
 \boxed{|K|\ge\frac{\pi r^2h}{3}
 =\frac1{18\pi\sqrt3}=0.01020979435966\ldots.}
\]
Every universal convex container must contain a rigid copy of this helix,
so the estimate applies to the universal-cover problem. The construction is classical in the study of convex hulls of space
curves; see Egerv\'ary~\cite{Egervary1950}.
No assertion that this helix maximizes convex-hull volume among all curves
is required or made here.

The chosen parameters are optimal for this lower-bound construction.
Putting $x=2\pi r$ gives $x^2+h^2=1$ and
$\pi r^2h/3=h(1-h^2)/(12\pi)$, whose maximum on $[0,1]$ is achieved at
$h=1/\sqrt3$.


This bound is weaker than the two-test-curve bounds in the main text. It is included to illustrate why a nonplanar curve alone forces positive volume and why convex-hull volume does not answer the universal-cover question.

\begin{figure}[htbp]
\centering
\includegraphics[width=.85\textwidth]{helix.pdf}
\caption{The helix lower bound uses tetrahedra with the two endpoints as a common vertical edge. Their projections are disjoint fan triangles in an inscribed polygon.}
\end{figure}
